\documentclass{article}
\usepackage{graphicx}
\usepackage{amsmath}
\usepackage{amssymb}
\usepackage{amsfonts}
\usepackage{hyperref}
\usepackage{exercises}
\usepackage{indentfirst}
\usepackage{amsthm}
\usepackage{csquotes}
\usepackage{yhmath}
\usepackage{comment}
\usepackage{placeins}
\usepackage{geometry}
\usepackage{tikz}
\newcommand*{\oldneg}{\mathord{\sim}}
\newcommand{\Z}{\mathbb{Z}}
\newcommand{\N}{\mathbb{N}}
\let\oldemptyset\emptyset

\theoremstyle{definition}
\newtheorem{definition}{Definition}

\setlength{\parindent}{.5in}
\geometry{margin=1in}
\setlength{\parskip}{12pt}
\title{Non-terms}
\begin{document}
\date{}
\maketitle

All literals are non-terms. For any non-terms $x$ and $y$, $xy,(x)y,x(y),(xy)$ are all non-terms.
The termified form of an expression $x_1x_2...x_{n-1}x_n$ where all $x_i$ are non-terms, is
$(mn)$ where $m$ is the termified form of $x_1x_2...x_{n-1}$ and $n$ is the termified form of 
$x_n$. The base cases for the termified forms are as follows: the termified form of $x$ is $x$ for all
literals, and, for all non-terms $x$, the termified form of $(x)$ is $n$ where $n$ is the 
termified form of $x$.

\end{document}
